\documentclass[11pt]{article}
\usepackage[T1]{fontenc}
\usepackage[utf8]{inputenc}
\usepackage{lmodern,amsmath,amssymb,amsthm,mathtools}
\usepackage[margin=1in]{geometry}
\usepackage{xurl}
\usepackage[hidelinks]{hyperref}
\usepackage{microtype}
\hypersetup{pdftitle={Where the quantum capacity sits in the arithmetical hierarchy},pdfauthor={Seth Douglas and Nidhal Mghirbi}}
\newcommand{\C}{\mathbb C}
\newcommand{\N}{\mathbb N}
\newcommand{\Z}{\mathbb Z}
\newcommand{\Q}{\mathbb Q}
\newcommand{\R}{\mathbb R}
\newcommand{\T}{\mathbb T}
\newcommand{\un}{\mathbf 1}
\newcommand{\tr}{\operatorname{tr}}
\newcommand{\ad}{\operatorname{ad}}
\newcommand{\dist}{\operatorname{dist}}
\newcommand{\supp}{\operatorname{supp}}
\newcommand{\Lip}{\operatorname{Lip}}
\newcommand{\Ga}{G_\alpha}
\newcommand{\au}{\alpha_u^+}
\newcommand{\aunif}{\alpha_u^{\mathrm{unif}}}
\newcommand{\eps}{\varepsilon}
\newcommand{\om}{\omega}
\newcommand{\omt}{\tilde\omega}
\newcommand{\omh}{\hat\omega}
\newcommand{\FU}{F_{\mathrm U}}
\newtheorem{theorem}{Theorem}[section]
\newtheorem{lemma}[theorem]{Lemma}
\newtheorem{proposition}[theorem]{Proposition}
\newtheorem{corollary}[theorem]{Corollary}
\newtheorem{conjecture}[theorem]{Conjecture}
\newtheorem{maintheorem}{Theorem}
\renewcommand{\themaintheorem}{\Alph{maintheorem}}
\newtheorem{maincorollary}[maintheorem]{Corollary}
\theoremstyle{definition}
\newtheorem{definition}[theorem]{Definition}
\newtheorem{remark}[theorem]{Remark}
\newtheorem{example}[theorem]{Example}
\setlength{\emergencystretch}{2em}
\raggedbottom
\newcommand{\RE}{\mathrm{RE}}
\newcommand{\coRE}{\mathrm{coRE}}
\title{Where the quantum capacity sits\\ in the arithmetical hierarchy}
\author{Seth Douglas \and Nidhal Mghirbi}
\date{Version 1.0.1, November 2026}
\begin{document}
\maketitle
\begin{abstract}
\input{abstract.txt}
\end{abstract}
\noindent{\footnotesize\textit{AI disclosure.} This paper was developed with AI assistance. The authors take responsibility for its claims; methods and verification limits appear in Section~\ref{sec:methods}.}
{\small\tableofcontents}
\section{Introduction}\label{sec:intro}
Can a finite description of a quantum channel tell us whether its quantum
capacity exceeds a threshold? There are two different questions here. An
algorithm may approximate a capacity from below without ever being able to
decide equality with the threshold. The distinction determines where the
decision problem belongs in the arithmetical hierarchy.

The picture is small. Write $\Sigma^0_1=\RE$ for the computably enumerable
sets, and $\Pi^0_1=\coRE$ for their complements. A predicate in $\Pi^0_2$
has the form $\forall k\,\exists s\,R(k,s)$ with $R$ decidable. The following
table gives upper bounds, except where completeness is expressly stated.
The ordinary capacity $Q$ is measured in qubits per channel use.
\begin{center}
\small
\begin{tabular}{p{0.62\linewidth}l}
\hline
Question & Class or status\\ \hline
$Q(\Phi)>0$; $Q(\Phi)=0$ & $\Sigma^0_1$; $\Pi^0_1$\\
$Q(\Phi)\geq\tfrac12$; $Q(\Phi)<\tfrac12$ & $\Pi^0_2$; $\Sigma^0_2$\\
$Q(\Phi)\geq\tfrac34$ versus $Q(\Phi)\leq\tfrac14$ (promise) & $\RE$~\cite{BMZ}\\
Closed thresholds for $C_{0E}(G)$ and $\Theta_q(G)$ & $\Pi^0_2$\\
One-shot game-graph ceiling, $q$ / $qc$ models & $\RE$-complete / $\coRE$-complete\\
$Q(\Phi_L)=0$ for every MPO-family index $L$ & $\coRE$, undecidable\\
\hline
\end{tabular}
\end{center}

\paragraph{Input and reduction conventions.}
All dimensions are finite and supplied. For ordinary capacity we allow
\emph{effective channel names}: computable Kraus entries with programs
giving certified approximations at requested precision, under the promise
that they describe a CPTP map. Computable Choi names are equivalent for
this purpose. Our default discrete domain $\mathcal E$ consists of explicit
algebraic Choi matrices, encoded by integer polynomials and isolating data;
validity is decidable. Circuit descriptions over a fixed computable gate
set are also covered; Clifford+$T$ gives algebraic data. Zero-error results
use explicit finite graphs or exact algebraic channel data. Hardness means
computable many-one hardness unless stated otherwise. Promise reductions
must produce valid instances; Turing reductions must make valid queries.
These conventions apply throughout.

\paragraph{Main contribution: the channel application and its hierarchy consequence.}
We apply standard effective finite-dimensional optimization and the
effective Fekete principle to quantum capacity. This gives increasing
rational lower bounds uniformly from a channel description and yields
the second-level closed-threshold bound below.
In particular, the problem called (QCap) by Bhattacharyya,
Mehta and Zhao (BMZ) lies in $\Pi^0_2$. Their Conjecture~1.4 reads
\begin{quote}
``(QCap) is $\Pi^0_3$-hard.''
\end{quote}
For their effective circuit setting, this is false as printed, under
computable many-one or Turing reductions. The correction concerns the
hierarchy discussion, rather than their hardness constructions.

\begin{remark}[Known ingredients and closest prior approximation result]
The coherent-information formula is due to Lloyd, Shor and
Devetak~\cite{Lloyd,Shor,Devetak}. Our finite-dimensional optimization
uses Winter's entropy continuity bound~\cite{Winter}.
Wolf--Cubitt--P\'erez-Garc\'ia~\cite{WCPG} distinguish exact decisions from
approximation. For quantum capacity, Khabbazi Oskouei--Mancini
\cite[Theorem~4.1]{OM} establish recursive approximability as convergence
of a computable rational sequence without a modulus or monotonicity.
The generic effective transformation from computable superadditive
representations to lower approximations is explicit in Boche, B\"ock and
Deppe~\cite[Lemmas~17--19 and Theorem~21]{BBD}, building on Zheng and
Weihrauch. Our contribution here is its channel-specific implementation
and the resulting correction to the QCap hierarchy discussion.
\end{remark}

\paragraph{New synthesis: what an undecidability proof would require.}
Positivity is semidecidable. On $\mathcal E$, it is decidable exactly when
zero capacity has a sound and complete effective system of finite
certificates. A separate activation-bound characterization has one
unconditional direction and one direction conditional on decidability of
the real exponential field. An undecidable all-length predicate already
exists for precisely specified MPO-indexed channel families. Its external
index is essential to the statement.

\begin{remark}[Known in special settings; recoveries]
BMZ prove QMA-hardness of a gapped ordinary-capacity problem and RE-hardness
of a restricted one-shot zero-error problem~\cite{BMZ}.
Cubitt--Elkouss--Matthews--Ozols--P\'erez-Garc\'ia--Strelchuk
\cite{CEMOPS} prove that no fixed block length suffices across all channels;
this does not give one channel that fails to attain capacity at every
length. The one-shot game-graph equivalences and model separations below
recover Harris~\cite{Harris}, Kim--Paulsen--Schafhauser~\cite{KPS}, and
Paddock--Slofstra~\cite{PS}. Our role in that part is the communication
phrasing.
\end{remark}

The third takeaway is a concrete obstruction to extending one-shot
hardness to asymptotic zero-error coding. A sufficient OR-repetition
conjecture remains unproved. Its classical shadow reaches a specific
24-vertex magic-square graph: all the evaluated spectral families give
the upper bound $6$, while a certified computation gives strong-square
independence number $29$. The resulting dichotomy isolates the missing
spectral information without deciding the Shannon capacity.

\section{Uniform lower approximation and closed thresholds}\label{sec:hierarchy}
For a channel $\Phi:A\to B$, choose a complementary channel $\Phi^c$ and put
\[
 b_n(\Phi)=\max_{\rho} \{S(\Phi^{\otimes n}(\rho))
                  -S((\Phi^c)^{\otimes n}(\rho))\},\qquad
 a_n(\Phi)=b_n(\Phi)/n.
\]
All logarithms are base two unless $\ln$ is written. Pure inputs give
$b_n\geq0$. Product inputs give $b_{n+m}\geq b_n+b_m$, so
$b_n$ is nondecreasing, although $a_n$ need not be. The capacity formula
and Fekete's lemma give
\begin{equation}\label{eq:capacity}
 Q(\Phi)=\lim_{n\to\infty}a_n(\Phi)=\sup_{n\geq1}a_n(\Phi).
\end{equation}

\begin{theorem}[Uniform finite-block and capacity approximation]\label{thm:uniform}
There is a uniform algorithm returning rationals $c(\Phi,n,k)$ such that
\begin{equation}\label{eq:c}
 |c(\Phi,n,k)-a_n(\Phi)|\leq2^{-k}.
\end{equation}
Consequently $Q$ is uniformly left-c.e.: a computable increasing rational
sequence $r_s(\Phi)$ converges to $Q(\Phi)$.
\end{theorem}
\begin{proof}
Here is the finite-block algorithm; Appendix~\ref{app:effective} supplies
two explicit effective-net derivations. A reference system of dimension
$D=d_A^n$ purifies every input. Thus $b_n$ is the maximum of
$-S(R|B)$ over a unit sphere in $\C^D\otimes\C^D$.
For unit vectors at Euclidean distance at most $\eta\leq1$, their output
states have trace distance at most $\eta$. Winter's bound gives the
computable modulus
\begin{equation}\label{eq:modulus}
 |I_c(\psi)-I_c(\varphi)|\leq
 2\eta n\log d_A+g(\eta),\quad
 g(\eta)=(1+\eta)h_2\!\left(\frac{\eta}{1+\eta}\right),
\end{equation}
where $g(0)=0$. Optimize certified entropy evaluations over a sufficiently
fine finite net. Both the mesh error and evaluation error are controlled
uniformly. Divide by $n$ to obtain~\eqref{eq:c}.

The following lower-approximation step is the effective Fekete principle
\cite[Lemmas~17--19 and Theorem~21]{BBD}, applied uniformly to the
finite-block channel quantities. Explicitly, define
\begin{equation}\label{eq:rs}
 r_s=\max\bigl(\{0\}\cup
 \{c(\Phi,n,k)-2^{-k}:1\leq n,k\leq s\}\bigr).
\end{equation}
These are increasing rational lower bounds. Given any $a_n$, taking
$k$ arbitrarily large shows $\sup_s r_s\geq a_n$.
Equation~\eqref{eq:capacity} proves convergence to $Q$.
\end{proof}

\begin{corollary}[Threshold bookkeeping]\label{cor:threshold}
For every fixed computable threshold $t$, the predicate $Q(\Phi)>t$ is
c.e., $Q(\Phi)\geq t$ is in $\Pi^0_2$, and $Q(\Phi)<t$ is in
$\Sigma^0_2$. In particular, (QCap) is in $\Pi^0_2$,
(GapQcap) is in RE, and $Q=0$ is co-RE.
\end{corollary}
\begin{proof}
For rational $t$ the decisive formulas have decidable matrices:
\begin{align}
 Q>t&\ \Longleftrightarrow\ \exists s\ [r_s>t],\label{eq:strict}\\
 Q\geq t&\ \Longleftrightarrow\
 \forall k\geq1\ \exists s\ [r_s>t-2^{-k}],\label{eq:closed}\\
 Q<t&\ \Longleftrightarrow\
 \exists k\geq1\ \forall s\ [r_s\leq t-2^{-k}].\label{eq:below}
\end{align}
For computable $t$, strict comparison with it is c.e.; its precision
witness merges with the existential stage in~\eqref{eq:closed}.
Complementation gives the last threshold class. Search for $r_s>1/2$
accepts every promised $Q\geq3/4$ instance and no promised $Q\leq1/4$
instance. Finally $Q\geq0$, so~\eqref{eq:strict} at zero proves the
positivity and zero-capacity claims.
\end{proof}
The RE upper bound for (GapQcap) is already stated by
BMZ~\cite[Section~1.3]{BMZ}; we include it to compare the gapped and
closed thresholds.

Equality at $Q=t$ is included in~\eqref{eq:closed}: every threshold
$t-2^{-k}$ is eventually crossed. This is why numerical approximation
does not turn the closed threshold into an existential test. For general
program names one can dovetail completed lower evaluations, including
partial computations off the valid-input promise. This gives a
$\Pi^0_2$ extension agreeing on valid names. No validity test for such
names is required.

\paragraph{BMZ's conjecture.}
Computable preimages preserve $\Pi^0_2$. Thus $\Pi^0_3$ many-one hardness
of (QCap) would collapse the strict arithmetical hierarchy. Under Turing
reductions, a $\Pi^0_2$ language is computable in $\mathbf0''$, whereas
a $\Pi^0_3$-complete language has degree $\mathbf0'''$. The same
contradiction holds for valid-query promise reductions, using the
extension just described. BMZ's stated $\Pi^0_4$ membership remains a
valid loose bound. Their separate gap-hardness conjecture is unaffected.

\begin{conjecture}[A corrected form]\label{conj:qcap}
On the effective algebraic circuit domain, (QCap) is $\Pi^0_2$-complete.
\end{conjecture}
This is a proposed replacement, not a lower bound proved here.

\paragraph{Question 2 and finite attainment.}
BMZ ask whether a channel exists with $a_n(\Phi)<Q(\Phi)$ for every
integer $n\geq1$. The following implication uses those integers, rather
than only powers of two.
\begin{proposition}\label{prop:attainment}
If every channel in the encoded family attains $Q$ at some finite block
length, then (QCap) is in $\Delta^0_2$. Hence $\Pi^0_2$ hardness, under
many-one or Turing reductions, would force a channel in that family
with $a_n<Q$ for all $n$.
\end{proposition}
\begin{proof}
Under finite attainment, $Q\geq t$ is equivalent to
$\exists n\ [a_n\geq t]$. For rational $t$,
\[
 a_n\geq t\ \Longleftrightarrow\
 \forall k\ [c(\Phi,n,k)+2^{-k}\geq t].
\]
Thus the threshold is $\Sigma^0_2$ and, independently by
Corollary~\ref{cor:threshold}, $\Pi^0_2$. Their intersection is
$\Delta^0_2$, computable in $\mathbf0'$. This excludes $\Pi^0_2$
Turing hardness as well as many-one hardness.
\end{proof}
No all-length nonattaining channel follows without the hardness
hypothesis. Nor does undecidability of an exact threshold alone imply
noncomputability of the capacity function: equality tests for computable
reals can already be undecidable. By contrast, many-one RE-hardness of
this closed threshold would exclude a uniformly computable $Q$, since
its closed threshold would then be co-c.e.; RE-hardness of the gapped
promise would also exclude it, since a computable $Q$ decides that promise.

\section{Zero certificates and activation}\label{sec:positive}
For the rest of this section the discrete domain is $\mathcal E$.
Define the activation length
\[
 n_*(\Phi)=\min\{n\geq1:b_n(\Phi)>0\},
\]
with value $\infty$ when $Q=0$. Positivity has a finite witness, although
the witness length is not supplied by the channel description.

\begin{proposition}[Zero-certificate equivalence; unconditional]\label{prop:zero}
The predicate $Q>0$ is decidable if and only if
$Z=\{\Phi:Q(\Phi)=0\}$ is c.e. Equivalently, zero capacity has a sound
and complete system of finite certificates with a computable verifier.
\end{proposition}
\begin{proof}
The complement of $Z$ is c.e. by Corollary~\ref{cor:threshold}.
If $Z$ is c.e., dovetail its enumeration with the positive-capacity
search. Exactly one terminates. Conversely a decision algorithm enumerates
$Z$. A c.e. enumeration yields finite certificates by recording the
enumeration stage, and a computable verifier yields an enumeration by
searching all finite certificates.
\end{proof}

\begin{remark}[Known structural certificates and a recent claim]
PPT and antidegradability give effective zero-capacity certificates.
For algebraic data, PPT is an exact positive-semidefiniteness test on the
partially transposed Choi matrix. Antidegradability is an exact
semialgebraic feasibility test for a symmetric two-output extension of
the normalized Choi state. Both are decidable, by real quantifier
elimination rather than approximate SDP termination~\cite{WCPG}.

Zhu--Wang~\cite[Theorems~1.1 and 4.3; Section~5]{ZW} claim that the
rational qutrit channel
\[
 \Lambda(X)=\frac12X+\frac14\bigl(\tr(X)I-X^{\mathsf T}\bigr)
\]
has zero private and quantum capacities while being NPT and
non-antidegradable. Their proposed certificate uses a signed lift,
a completely positive defect factorization, and a complete information
comparison that tensorizes. We cite this as a preprint claim; its full
signed-lift/BKM argument is not independently verified here and is not
an input to our proofs. It changes the appropriate structural question:
the issue for Proposition~\ref{prop:zero} is an effective complete
enumeration of all zero-capacity channels. Additional examples or
certificate classes can enlarge that enumeration without making it
complete.

Sandhu's manuscripts~\cite{SandhuOrder,SandhuCapacity}, which state that
their results were developed by AI systems, claim more: complete channel
order is undecidable, and on finite rational-complex channel descriptions
$Q>0$ and $P>0$ are $\RE$-complete, with the matching statements for
one-way distillable entanglement and key. If this holds, the enumeration
in Proposition~\ref{prop:zero} does not exist on $\mathcal E$, which
contains those descriptions: it is the reduction to positivity that the
synthesis in Section~\ref{sec:intro} describes. A refutation-first check
of ours found no error in the chain from complete order to exact zero
capacity; the separate fixed-error inapproximability theorem
of~\cite{SandhuOrder} was not checked. Neither claim is an input to our
proofs. Hirche~\cite{Hirche} relates complete entropic channel orders to
binary estimation and studies zero private capacity, activation, and
sharp antidegradability thresholds for Werner--Holevo channels.
\end{remark}

The obstruction can be stated more sharply without assuming any
particular structural classification.
\begin{proposition}[A fixed certificate family]\label{prop:escape}
Suppose $e\mapsto\Phi_e$ is a computable many-one reduction from HALT to
$\{Q>0\}$. For every fixed sound c.e. set $C\subseteq Z$, the set
\[
 D_C=\{e:e\notin\mathrm{HALT},\ \Phi_e\notin C\}
\]
is not c.e., and is infinite. If $C$ includes the PPT and antidegradable
channels, infinitely many distinct channel outputs evade both classes.
\end{proposition}
\begin{proof}
The preimage of $C$ is c.e. and consists of nonhalting indices. Its union
with $D_C$ is the complement of HALT, which is not c.e. If only finitely
many distinct outputs occurred in $D_C$, exact equality of algebraic
channel matrices with those finitely many outputs would enumerate $D_C$,
a contradiction.
\end{proof}
This is a statement about uniform completeness. It does not prohibit a
proof of zero capacity for an individual channel. The elementary
certificate principle is analogous to the distillability discussion in
\cite[Section~3.1]{WCPG} and the tensor-stable-positivity discussion in
\cite[Section~4]{VNC}.

Let $|\Phi|$ denote description length and set
\[
 A(s)=\max\{n_*(\Phi):|\Phi|\leq s,\ Q(\Phi)>0\},
\]
with maximum zero when the set is empty. Finitely many descriptions have
length at most $s$, so $A(s)$ is a finite number.

\begin{proposition}[Activation bounds; conditional reverse direction]\label{prop:activation}
If positivity is decidable, $A$ is bounded by a computable function.
Conversely, assume that $\operatorname{Th}(\R, +,\cdot,<,\exp)$ is
decidable. Then a computable bound on $A$ decides positivity. In
particular the equivalence holds under Schanuel's conjecture, by the
Macintyre--Wilkie theorem.
\end{proposition}
\begin{proof}
With a positivity decider, list all valid descriptions of length at most
$s$ and identify the positive ones. Run the lower-bound searches for
those finitely many channels until each has a witness $b_n>0$.
The largest witnessed length, or zero for an empty list, is a computable
upper bound on $A(s)$.

For the converse, fixed-block positivity is an existential sentence over
the real exponential field. Purify the input in dimension $D\times D$.
Its output matrices are polynomial in the real and imaginary input
coordinates, with algebraic coefficients. Quantify a unitary
diagonalization and nonnegative eigenvalues $\lambda_i$ for each matrix.
Represent $\lambda\ln\lambda$ by $t=\lambda y$, $\exp(y)=\lambda$ when
$\lambda>0$, and by $t=0$ when $\lambda=0$.
Algebraic coefficients have finite isolating definitions. The resulting
sentence expresses positive coherent information; changing the logarithm
base does not change its sign. The assumed theory decider therefore
decides $b_n>0$. If $A\leq f$ computably, test
$b_{\max(1,f(|\Phi|))}>0$. Monotonicity of $b_n$ proves correctness.
\end{proof}
\begin{remark}[Conditionality and attribution]
Decidability of the real exponential field is not supplied by
semialgebraic feasibility. Macintyre--Wilkie's implication from Schanuel
is conditional~\cite{MW}; see also~\cite[version~3, Theorem~14.4]{BG}.
The proposition is a bounded-witness consequence with a channel-specific
entropy encoding. Under that hypothesis, undecidability of positivity
would mean that $A$ is not eventually bounded by any computable function:
finitely many exceptional sizes could be hardcoded. It does not say that
every positive channel requires long activation.
\end{remark}

\section{An external index: MPO-generated channel families}\label{sec:mpo}
An all-length question about a family can be undecidable even when every
zero-capacity member has a structural certificate. The following wrapper
makes the distinction explicit.

\begin{proposition}[Measure-and-erase wrapper]\label{prop:wrapper}
Let $W$ be a finite Hermitian matrix and let $c>0$ be rational with
$c\geq\|W\|$. There is a channel $\Phi[W]$ with algebraic data when $W$
is algebraic, and rational Choi data when $W$ is rational, such that
\[
 Q(\Phi[W])=0\ \Longleftrightarrow\ W\geq0.
\]
\end{proposition}
\begin{proof}
Put $A=(I-W/c)/2$ and $R=(I+W/c)/2$. On an input $X\otimes q$, with
$q$ a message qubit, measure the effects $A,R$ on $X$. Bob receives
the qubit and an accept flag in the first branch, and an erasure and
reject flag in the second. An isometry sends the measured $X$ to the
environment, with $\sqrt A$ or $\sqrt R$ applied; the environment gets
the qubit exactly in the reject branch. Orthogonal flags in both outputs
remove cross terms. Bob's map is
\[
 \rho\longmapsto
 |\mathrm{acc}\rangle\langle\mathrm{acc}|\otimes
       \tr_X[(A\otimes I)\rho]
 +\tr[(R\otimes I)\rho]\,
       |\mathrm{rej},e\rangle\langle\mathrm{rej},e|.
\]
It is linear in $A,R$, which proves the rational Choi assertion.

If $W\geq0$, then $A\leq R$ and $R\geq I/2$. On its reject branch
the environment uses the effect $G=R^{-1/2}AR^{-1/2}\leq I$ on $X$,
forwards its qubit on success, and erases otherwise. On its accept
branch it erases. The identity $\sqrt R G\sqrt R=A$ recovers Bob's
accept map, including correlations with the message; the remaining
weight is his reject map. This is an antidegrading channel.

If $\langle v,Wv\rangle=-\mu<0$ for a unit vector $v$, fix $X=v$.
The message sees a qubit erasure channel with acceptance probability
$(1+\mu/c)/2$. The maximally mixed message input has coherent information
$\mu/c>0$.
\end{proof}

\begin{corollary}[Indexed-family undecidability]\label{cor:mpo}
Given a rational TI MPO description generating diagonal Hermitian
operators $W_L$, set $\Phi_L=\Phi[W_L]$ with effective positive rational
norm bounds $c_L$. The predicate $\forall L\geq1\ Q(\Phi_L)=0$ is co-c.e.
and undecidable. For rational open-boundary descriptions the complementary
existential predicate is RE-complete. Periodic fixed-size descriptions
also give undecidability, without a many-one-completeness claim here.
\end{corollary}
\begin{proof}
Kliesch--Gross--Eisert~\cite[Theorem~5 and Lemmas~11--12]{KGE} give a
many-one reduction to existence of a negative eigenvalue at some length for rational
open-boundary diagonal MPOs, even at physical dimension two and fixed
bond bound. Finite-length positivity is decidable exactly, so existence
of a negative eigenvalue is c.e. and RE-complete. The wrapper preserves
this predicate length by length.

De las Cuevas et al.~\cite[Problem~2 and Lemma~3]{DCCWP} give a periodic
version with seven rational $7\times7$ matrices. Their promise separates
positivity at every length from a negative eigenvalue at every sufficiently large
length. Their reduction framework permits Turing reductions; this
supports the periodic undecidability assertion. If the matrices are
$C_i\in\Q^{s\times s}$, take
$K=\max(1,\max_i\|C_i\|_\infty)$ and $c_L=sK^L$ for trace contraction.
For open boundaries take
$c_L=\max(1,\|\ell\|_1\|r\|_\infty)K^L$. These bound every diagonal
coefficient. Complementation proves co-c.e. membership of the universal
predicate.
\end{proof}

The translation invariance is that of the generating MPO tensor.
The channel has a distinguished message qubit and terminal flags; its
Choi matrix has fixed-bond MPO form on the $X$ chain with terminal
auxiliary systems. This specifies an effectively generated family,
rather than a spatially translation-invariant local channel with
identical input and output sites. In particular $\Phi_L$ is not asserted
to be $\Phi_1^{\otimes L}$. Every universal-YES length is antidegradable;
undecidability remains in the external quantifier over $L$.

\begin{remark}[Attribution]
The undecidability is imported from MPO positivity. Transport of that
problem to tensor positivity already appears in~\cite{VNC}. The wrapper
adds a CPTP capacity interpretation. The memory-channel capacity theorem
of Elkouss--P\'erez-Garc\'ia~\cite{EPG} is an operational analogue, for
classical capacity with memory.
\end{remark}

\section{Zero error: finite witnesses and the missing repetition step}\label{sec:zeroerror}
For a finite confusability graph $G$, write $\alpha_*(G)$ for its
one-use message maximum with arbitrary finite-dimensional entanglement
assistance. Write $\alpha_q(G)$ for the restricted maximally entangled,
projective code parameter. These are different resources.
Beigi's entanglement-assisted theta bound~\cite{Beigi} also bounds the
restricted parameter. For recent resource-sensitive game-to-classical-channel
correspondences and an arbitrary-versus-maximally-entangled separation,
see Man\v{c}inska~\cite[Theorem~6 and Corollary~8]{Mancinska}; we do not
use that result in the proofs below. Set
\[
 C_{0E}(G)=\sup_m\frac{\log\alpha_*(G^{\boxtimes m})}{m},\qquad
 \Theta_q(G)=\sup_m\alpha_q(G^{\boxtimes m})^{1/m}.
\]
The strong product has distinct tuples adjacent when every coordinate
is equal or adjacent. Product protocols give the required
supermultiplicativity.

\begin{proposition}[Uniform zero-error lower approximation]\label{prop:zerolower}
On explicit finite graphs, $C_{0E}$ and $\Theta_q$ are uniformly left-c.e.
Their closed computable-threshold predicates lie in $\Pi^0_2$.
The same lower-approximation conclusion for $C_{0E}$ holds with exact
algebraic channel data.
\end{proposition}
\begin{proof}
For fixed message number $M$ and dimension $d$, entanglement-assisted
graph coding is finite real-algebraic feasibility. One may use positive
semidefinite matrices $\rho,\rho_i^v$ with
\[
 \tr\rho=1,\quad \sum_v\rho_i^v=\rho,\quad
 \rho_i^v\rho_j^w=0\quad
 (i\ne j,\ v=w\text{ or }v\sim w).
\]
For $\alpha_q$, replace these by projections $E_{i,v}$ with
$\sum_vE_{i,v}=I$ and the same cross-message zero products.
Real quantifier elimination decides each fixed-dimensional system.
Dovetail over $m,M,d$ and enumerate rational lower estimates for
$\log M/m$, or $M^{1/m}$. For algebraic channel data the fixed-dimensional
protocol constraints likewise have exact algebraic coefficients.
The supremum argument of~\eqref{eq:closed} completes the proof.
\end{proof}
\begin{remark}[Known criterion; representation restriction]
These code descriptions are standard~\cite{CLMW,MR,LZ}.
An approximation program for arbitrary computable-real entries does
not decide exact support or zero products. The zero-error statement
therefore uses the exact inputs specified in Section~\ref{sec:intro}.
No such restriction was needed for ordinary capacity. Spectral duality
in~\cite{LZ} is a minimum over an abstract spectrum; it supplies no
effective convergent family of upper bounds. Li--Li's multiplicative
SDP bound~\cite{LL} concerns no-signalling assistance and supplies a
single bound. Neither it nor the spectral families in our related
paper~\cite{MD} supply the missing convergence theorem.
\end{remark}

\paragraph{One-shot ceilings: a recovery.}
Normalize a finite synchronous game by retaining self-legal answers,
rejecting unequal answers at equal questions even off the sampling
support, and making only unsupported off-diagonal tests vacuous.
If a question loses every self-answer, detect it and use a fixed NO
instance. Its graph $X$ has vertices $(q,a)$ and edges at distinct pairs
for which either orientation is forbidden. There are $N=|Q|$ question
cliques. Here $q$ means finite-dimensional strategies, $qa$ their
closure, and $qc$ commuting-operator strategies. Harris's
Theorem~5.5~\cite{Harris} identifies a perfect strategy
with $\alpha_t(X)=N$ in the relevant $q,qa,qc$ models.

In communication terms, aggregate code projections into
$P_{q,a}=\sum_iE_{i,(q,a)}$ and $F_q=\sum_aP_{q,a}\leq I$.
At the ceiling, $\sum_q\tau(F_q)=N$ forces $F_q=I$ for a faithful
trace, giving perfect game measurements. Conversely the code is
$E_{i,(q,a)}=\mathbf1_{i=q}P_{q,a}$, also in the tracial commuting model.

The finite-dimensional ceiling is RE-complete; the commuting ceiling
is co-RE-complete, by the known satisfiability results
\cite{JNVWY,PS}; for the latter see~\cite[Theorem~5.11(c)]{PS}.
KPS~\cite[Theorem~5.4 and Corollary~5.5]{KPS}
give earlier $q$ versus $qa$ graph separations. Harris's
Corollary~5.6 gives a $qa$ versus $qc$ separation. These are recoveries,
not new separation or hardness theorems. Nonattainment of a finite-dimensional
perfect strategy at supremum value one does not meet the soundness
hypothesis below.

\paragraph{The precise missing step (conjectural).}
A projective packing consists of projections $P_v$ with $P_vP_w=0$
on graph edges, of weight $\sum_v\tau(P_v)$. Its supremum is
$\alpha_p$. Aggregating a projective code gives a packing, so
$\alpha_q\leq\alpha_p$. A packing does not supply the message-indexed
PVM decomposition needed for the converse.

\begin{conjecture}[Packing OR-repetition]\label{conj:orc}
For every normalized finite synchronous game with an effective rational
question distribution having uniform marginals and entangled value
$\omega^*\leq1/2$, there is $c_G\in(0,1)$ such that, for every $m\geq1$,
\[
 \alpha_p(X^{\boxtimes m})\leq(1-c_G)^m N^m.
\]
\end{conjecture}
The constant may depend on the game; no computable modulus is asserted.
The version with $\alpha_q$ would suffice for the restricted projective
capacity and may be weaker. The strong power corresponds to the
\emph{special OR predicate}: identical tuple vertices are accepted, as
are pairs with a coordinate that is distinct and nonadjacent. It is not
Boolean OR of coordinate acceptance. AND repetition instead gives the
disjunctive product, adjacent when some coordinate loses.

The verified line extraction gives success at least $2M/N^m-1$ from
a packing of weight $M$, after completing abstentions by fixed answers.
It uses $\tau(PQ)\geq\tau(P)+\tau(Q)-1$, without commutativity.
This is only a constant density loss. Anchored AND parallel
repetition~\cite{BVY} does not give exponential OR loss.
The attempted AND-enforcing route fails on product
inputs because tensor-product output overlaps multiply and hence give
the strong product.

JNVWY's Theorem~12.6~\cite{JNVWY} provides attained finite-dimensional
perfect completeness and soundness at most $1/2$ for the halting
reduction. Diagonal normalization preserves that witness. After
symmetrizing the rational distribution, a common marginal
$\nu(q)=k_q/K$ becomes uniform by replacing $q$ with $k_q$ clones.
Averaging clone-indexed measurements preserves the soundness bound.
Thus Conjecture~\ref{conj:orc} would separate the asymptotic $\Theta_q$
ceiling on this family. This is a conditional route, with no conclusion
here for unrestricted entanglement-assisted codes.

\section{The magic-square certificates}\label{sec:ms}
Let $X=X_{\rm MS}$. Its vertices are the four legal bit assignments to
each of $R_1,R_2,R_3,C_1,C_2,C_3$ in a $3\times3$ square; all context
parities are even except $C_3$. Distinct vertices are adjacent if they
belong to the same context or disagree on a shared cell. This fixes the
24-vertex synchronous graph, rather than another magic-square event
graph. Ordinary independence is $\alpha(X)=5$.

\begin{proposition}[All-field certificate]\label{prop:haemers}
For every field $\mathbb F$, the fractional Haemers bound satisfies
$\mathcal H_f^{\mathbb F}(X)=6$.
\end{proposition}
\begin{proof}
A $d$-representation is a block fitting matrix $M$ with diagonal $I_d$
and zero blocks at distinct nonneighbors; its cost is $\operatorname{rank}M/d$.
Factor $M=BA$ and put $V_v=\operatorname{im}A_v$. Then
$B_vA_v=I_d$ and $B_vA_w=0$ at nonneighbors. Consequently
\[
 r(S)=d^{-1}\dim\sum_{v\in S}V_v
\]
is normalized, monotone and submodular, with
\begin{equation}\label{eq:rankconstraints}
 r(\{v\})=1,\qquad
 r(\overline N(v)\cup\{v\})-r(\overline N(v))=1,
\end{equation}
where $\overline N(v)$ denotes distinct nonneighbors.

The induced subgraph on
\[
\begin{gathered}
 R_1:011;\quad R_2:000,011;\quad R_3:000,011;\quad C_1:000;\\
 C_2:101,110;\quad C_3:100,111
\end{gathered}
\]
is $2K_1\sqcup W_8$, where $W_8$ is the eight-cycle with opposite
chords. The dimension inequalities force cost at least $6$ on this core.
This is also the $n=8$ M\"obius-ladder Shannon-inequality argument of
Blasiak--Kleinberg--Lubetzky~\cite[Theorem~5.2 and proof]{BKL}, applied
to dimensions over any field. For an explicit finite verification, the
accompanying certificate gives the identity
\begin{equation}\label{eq:dual}
 r(T)-6=\sum_i\lambda_i g_i(r)+\sum_j\mu_j h_j(r),
 \qquad \lambda_i\geq0,
\end{equation}
using 33 monotonicity/submodularity rows $g_i\geq0$ and 10 rows
$h_j=0$ from~\eqref{eq:rankconstraints} and normalization.
Every coefficient is rational and the identity is checked exactly.
Restriction to the core gives the lower bound on $X$; its six question
cliques give the upper bound $6$.
\end{proof}
\begin{remark}[Recovery and exact quantum ceiling]
The subspace framework is standard~\cite[Section~2.1]{BC}; the
M\"obius-ladder lower bound is a recovery of~\cite{BKL}. The specific
induced-core application supplies a useful certificate, not a new LP
method. The Mermin--Peres observables~\cite{Mermin,Peres}, displayed in
Appendix~\ref{app:certs}, give an exact four-dimensional perfect
strategy. Therefore $\alpha_q(X)=\vartheta(X)=6$. Beigi's
entanglement-assisted theta bound~\cite{Beigi} also applies to this
restricted code parameter. Vertex transitivity
and clique number four give $\overline\chi_f(X)=24/4=6$.
These equalities use exact projections and combinatorial certificates,
rather than floating SDP estimates.
\end{remark}

\paragraph{Certified computation.}
The exact strong-square value is
\begin{equation}\label{eq:29}
 \alpha(X^{\boxtimes2})=29.
\end{equation}
A 29-word independent set was independently constructed in the first
blind check and verified pairwise. For the upper bound, slicing by
the first question gives at most $6\alpha(X)=30$ words. A question-pair
block contains at most one word. A 30-word code therefore has five
words in each row and column of question blocks, so its empty blocks
form a permutation. Checked graph symmetries reduce the 720 permutations
to two double cosets of sizes 72 and 648.
The earlier independent encodings were UNSAT with CaDiCaL and Glucose.
A further independent encoding was UNSAT with Lingeling for both
representatives, and drat-trim separately verified both saved proof
traces. The programs, formulas, traces, checker source, symmetry
verification and lower witness are included in the computational supplement.

\paragraph{Which spectral bounds are exhausted?}
Every evaluated known spectral family gives $6$ on $X$: $\vartheta$,
$\overline\chi_f$, all fractional Haemers bounds, and every point in
the quantum spectrum, the last squeezed between $\alpha_q(X)=6$ and
the question-clique ceiling. This also covers entanglement-assisted
spectral points, including those reported in~\cite{MD}, provided their
spectral membership is accepted. No proof in this note depends on that
preprint's new membership claim.

Vrana's refinements~\cite[Theorems~1.1--1.2]{Vrana} do not improve these
evaluations. For a vertex-transitive graph, concavity and averaging over
automorphisms show $F_f(X,U)=\log f(X)$ at the uniform distribution $U$.
Thus interpolation of refinements gives
$f_F(X)=\prod_i f_i(X)^{\lambda_i}=6$, also for pointwise limits of
these interpolations. This specifies the evaluated closure, without
claiming an exhaustive classification of all spectral constructions.

\begin{corollary}[Spectral dichotomy]\label{cor:dichotomy}
Either $\Theta(X)=6$, or a classical spectral point outside the evaluated
families has value below $6$ on $X$. In the second case
$\Theta(X)<\Theta_q(X)=6$. Presently the certified interval is
$\sqrt{29}\leq\Theta(X)\leq6$.
\end{corollary}
\begin{proof}
Zuiddam's duality~\cite{Zuiddam} says $\Theta(X)=\min_f f(X)$.
Equation~\eqref{eq:29} gives the lower endpoint; the preceding evaluations
give the assertions about a minimizing point.
\end{proof}
Vertex transitivity gives $\vartheta(\overline X)=24/\vartheta(X)=4$;
the clique of size four gives $\Theta(\overline X)=4$. The diagonal
independent set in $X\boxtimes\overline X$ has size 24, so the second
branch would also give
$\Theta(X\boxtimes\overline X)>\Theta(X)\Theta(\overline X)$.
Supermultiplicativity of this kind is known in general from Alon~\cite{Alon}.
The direct evaluated Haemers bounds fail to certify a strict gap here;
structural arguments using auxiliary graphs remain possible.

\paragraph{The classical shadow (conjectural).}
Qualitative classical OR-repetition for all full-support synchronous
games is equivalent to
\[
 \alpha(G)<\overline\chi(G)\quad\Longrightarrow\quad
 \Theta(G)<\overline\chi(G).
\]
A minimum clique cover can be partitioned into question cliques; a
perfect strategy is exactly an independent transversal. The only live
case is $\alpha<\overline\chi_f=\overline\chi$. This is the integral-cover
case of Shannon's $C_0$ versus $C_{0F}$ question~\cite{Shannon}, using
the classical channel whose output supports are all maximal cliques.
Then feedback fractional packing equals $\overline\chi_f$, and Shannon's
feedback formula applies when $C_0>0$. The realization matters: an
edge-symbol channel with the same graph can have a different feedback
packing value~\cite{CLMW}. For $X$, the two realizations give packing
values $6$ and $12$, respectively. This broad classical conjecture is
not equivalent to the narrower $\omega^*\leq1/2$ quantum conjecture;
the magic square itself is quantum-perfect. The first branch of
Corollary~\ref{cor:dichotomy} would refute the classical conjecture,
since $\alpha(X)=5<6$.

\section{Discussion}\label{sec:discussion}
The proved picture separates four issues. Ordinary capacity is
computably accessible from below; equality with a positive threshold
adds a universal precision quantifier. Positivity has finite witnesses,
but complete zero certificates need not be enumerable. Under the real
exponential hypothesis, computable bounds on activation would settle
positivity. In the MPO family, undecidability comes from an external
all-length quantifier. For restricted asymptotic zero-error coding, the
unproved issue is exponential packing loss in strong powers.

\paragraph{Outlook (conjectural and heuristic).}
An asymptotic projective zero-error ceiling hardness result remains open. Three developments would justify reopening the broader
search: a spectral point below six on $X_{\rm MS}$; a mechanism for
coherent amplification within tensor powers; or effective zero-capacity
certificates beyond PPT and antidegradability. The third condition must
now take the Zhu--Wang claim into account. An independently validated
new certificate could enlarge the structural toolkit; it would not by
itself establish completeness of that toolkit. The first condition
would settle the concrete classical test, not prove the quantum
repetition conjecture. The second is a missing mechanism, not an
existence theorem.

\emph{Heuristic.} Distillability, tensor-stable positivity, $Q>0$ and
quantum $C_0>0$ may be manifestations of a common problem about witnesses
in tensor powers. Their analogies suggest looking simultaneously for
positive finite witnesses and effective zero certificates. No reductions
or equivalence between these four questions are claimed. In particular,
incompleteness of a named structural certificate class does not alone
imply undecidability.

\section{Methods and provenance}\label{sec:methods}
\subsection*{How this work was done}
The authors developed the results and proofs with assistance from AI
models from more than one developer. They take responsibility for every
claim. The authors chose the question, directed the work and decided
which lines to pursue; the models did most of the exploration,
derivation and drafting.

Three rounds of proof development recorded arguments, unsuccessful approaches and self-checks.
Separate sessions that had not seen the material were instructed to
refute it and check primary sources. This version incorporates their
repairs. Literature and priority searches were AI-assisted and bounded
to the checked sources; recent preprint claims were not accepted as
independently verified proofs.
A subsequent independent whole-paper audit verified the finite certificates
and both strong-square UNSAT traces with a separate checker, and requested
the attribution and reference repairs recorded with that audit.

\paragraph{Computational evidence.}
The all-field dimension certificate uses only rational arithmetic after
LP discovery; its semantic replay reconstructs the graph and checks
every inequality row and coefficient of~\eqref{eq:dual}. This replay
and the exact non-diagonal wrapper check were run again during assembly.
The square upper bound has independent CaDiCaL, Glucose and Lingeling
encodings; a separate drat-trim check verifies both saved Lingeling
UNSAT traces. The supplement contains the finite certificate data,
checker programs and source, formulas, traces, symmetry verification,
lower witness and replay instructions. The graph reduction, encodings
and checker implementation remain part of the trust base; this is not
a proof-assistant formalization. These
computations certify the stated finite quantities; none extrapolates a
Shannon-capacity gap from finitely many powers.

\appendix
\section{Two effective finite-block derivations}\label{app:effective}
In the purification proof, grid the real unit sphere in dimension
$2D^2$. Choose rational mesh small enough that rounding a unit vector
has Euclidean error at most $\eta/2$, discard zero vectors, and normalize.
Normalization changes the error by at most another $\eta/2$.
The resulting finite algebraic net is effective. To evaluate entropy,
approximate each computable Hermitian output by a rational Hermitian
matrix in operator norm. Isolate its algebraic eigenvalues and use Weyl's
inequality to bound the ordered true eigenvalues. Intersect the
eigenvalue intervals with $[0,1]$. The continuous extension of $-x\log x$ at zero has an effective modulus on $[0,1]$;
zero eigenvalues cause no rank decision. Choose the net and evaluation
errors to sum to the requested tolerance in~\eqref{eq:modulus}.

The independent complementary-channel derivation starts with a
computable positive square root of the unnormalized Choi matrix
$J_\Phi=(\operatorname{id}\otimes\Phi)(|\Omega\rangle\langle\Omega|)$,
where $|\Omega\rangle=\sum_{i=1}^{d_A}|i\rangle|i\rangle$.
Reshape all $d_Ad_B$ square-root columns into Kraus operators, including
zero columns, to obtain a fixed-size environment $d_E=d_Ad_B$.
If starting from the normalized Choi state $J_\Phi/d_A$, multiply its
square-root columns by $\sqrt{d_A}$ before reshaping. No minimal-rank decomposition is
needed. For input states at trace distance $\delta\leq1$, separate
entropy continuity bounds on the $B$ and $E$ outputs give
\[
 |I_c(\rho)-I_c(\sigma)|
 \leq2n\delta\log(d_Bd_E)+2g(\delta).
\]
Density matrices have effective nets $YY^*$ with $Y$ drawn from
normalizations of nonzero rational complex matrices in Frobenius norm:
every state is $YY^*$ for
some unit-Frobenius $Y$, and
$\tfrac12\|YY^*-ZZ^*\|_1\leq\|Y-Z\|_F$.
Finite maximization with this modulus again proves~\eqref{eq:c}.
An even shorter lower-only proof enumerates dense algebraic inputs and
certified coherent-information lower evaluations over every length.
Continuity identifies their supremum with $Q$.

\section{Certificate conventions}\label{app:certs}
The exact magic-square observables, with Pauli matrices $X,Y,Z$, are
\[
 \begin{pmatrix}
 X\otimes I&I\otimes X&X\otimes X\\
 I\otimes Z&Z\otimes I&Z\otimes Z\\
 X\otimes Z&Z\otimes X&Y\otimes Y
 \end{pmatrix}.
\]
They commute within contexts and multiply to $I$ except in the last
column, where the product is $-I$. Their legal joint eigenspace
projections have rank one and give the asserted perfect code.

For the all-field certificate, subsets of the ten core vertices are
encoded by bit masks in the order printed in Section~\ref{sec:ms}.
The accompanying \texttt{dimension-dual.json} records all nonzero
multipliers and rows. Monotonicity rows are
$r(S\cup\{i\})-r(S)$; submodularity rows are
$r(S\cup\{i\})+r(S\cup\{j\})-r(S)-r(S\cup\{i,j\})$.
The replay verifies that the equality rows correspond to the graph's
actual nonneighbors before summing them. This gives a finite exact
proof of the numerical LP optimum, independent of floating solver
accuracy. The computational supplement includes the SAT formulas,
symmetry builder, pairwise lower-witness check, both saved Lingeling
UNSAT traces and drat-trim checker source. The two traces were separately
verified by drat-trim. This reduces reliance on solver outcomes, while
the graph reduction, Boolean encodings and checker implementation remain
in the trust base; no proof-assistant verification is claimed.

\bibliographystyle{plainurl}
{\small\bibliography{references}}
\end{document}
